\documentclass[11pt]{article}
\usepackage[margin=0.9in]{geometry}
\usepackage{amsmath,amssymb,amsthm,mathtools}
\usepackage[T1]{fontenc}
\usepackage[utf8]{inputenc}
\usepackage{enumitem}
\usepackage{booktabs,longtable,array,seqsplit}
\usepackage{hyperref}
\usepackage{xcolor}
\usepackage{microtype}
\hypersetup{colorlinks=true,linkcolor=blue!45!black,citecolor=blue!45!black,urlcolor=blue!45!black}
\setlength{\emergencystretch}{6em}

\newtheorem{theorem}{Theorem}[section]
\newtheorem{proposition}[theorem]{Proposition}
\newtheorem{lemma}[theorem]{Lemma}
\newtheorem{corollary}[theorem]{Corollary}
\newtheorem{claim}[theorem]{Claim}
\theoremstyle{definition}
\newtheorem{definition}[theorem]{Definition}
\newtheorem{assumption}[theorem]{Assumption}
\theoremstyle{remark}
\newtheorem{remark}[theorem]{Remark}

\newcommand{\R}{\mathbb{R}}
\newcommand{\C}{\mathbb{C}}
\newcommand{\N}{\mathbb{N}}
\newcommand{\Hsp}{\mathcal{H}}
\newcommand{\Fsp}{\mathcal{F}}
\newcommand{\Id}{\mathrm{Id}}
\newcommand{\tr}{\operatorname{tr}}
\newcommand{\diag}{\operatorname{diag}}
\newcommand{\rank}{\operatorname{rank}}
\newcommand{\spec}{\operatorname{spec}}
\newcommand{\spanop}{\operatorname{span}}
\newcommand{\dist}{\operatorname{dist}}
\newcommand{\range}{\operatorname{ran}}
\newcommand{\Ker}{\operatorname{ker}}
\newcommand{\Pinch}{\mathcal{C}}
\newcommand{\Off}{\widetilde{\mathcal{C}}}
\newcommand{\Pperp}{P^{\perp}}
\newcommand{\Qperp}{Q^{\perp}}
\newcommand{\op}{\mathrm{op}}
\newcommand{\F}{\mathrm{F}}
\newcommand{\HS}{\mathrm{HS}}
\newcommand{\ip}[2]{\left\langle #1,#2\right\rangle}
\newcommand{\norm}[1]{\left\lVert #1\right\rVert}
\newcommand{\abs}[1]{\left\lvert #1\right\rvert}
\newcommand{\code}[1]{\nolinkurl{#1}}

\newenvironment{argument}{\paragraph{Argument.}\begin{enumerate}[leftmargin=2em]}{\end{enumerate}}
\newenvironment{proofoutline}{\paragraph{Proof architecture.}\begin{enumerate}[leftmargin=2em]}{\end{enumerate}}
\newcommand{\sourceanchor}[1]{\par\smallskip\noindent\textbf{Source anchor.} #1\par\smallskip}
\newcommand{\sourcecomment}[1]{\begin{remark}#1\end{remark}}
\newenvironment{sourceguide}{\begin{description}[style=nextline,leftmargin=3.3cm,labelwidth=3.1cm]}{\end{description}}
\newenvironment{formalizationmap}{\begin{longtable}{@{}p{0.20\textwidth}p{0.31\textwidth}p{0.40\textwidth}@{}}\toprule Source item & Modern mathematical content & Repository status \\\midrule\endfirsthead\toprule Source item & Modern mathematical content & Repository status \\\midrule\endhead}{\bottomrule\end{longtable}}

\title{Double-angle Davis--Kahan theory:\\reflection defects, generic separation, tangent branches, and canonical spectral selection}
\author{}
\date{Comparative formalization dossier}
\begin{document}
\maketitle

\section*{Sources and purpose}
\begin{sourceguide}
\item[Primary source] Davis--Kahan (1970), especially the Section 7 reflection proof and Theorems 8.1--8.2.
\item[Generic separation] Albrecht Seelmann, \emph{Notes on the $\sin 2\Theta$ Theorem}, arXiv:1310.2036.
\item[Off-diagonal tangent] L. Grubi\v{s}i\'c, V. Kostrykin, K. A. Makarov, and K. Veseli\'c, \emph{The $\tan 2\Theta$ Theorem for Indefinite Quadratic Forms}, arXiv:1006.3190.
\item[Companion literature] Albeverio--Motovilov on generic scalar $\sin 2\theta$ bounds and a priori tangent theorems.
\item[Purpose] Determine the valid parent statements for $\sin 2\Theta$ and $\tan 2\Theta$, explain why a residual theorem needs a Ritz/Galerkin structure, and isolate the continuation theorem that selects the canonical branch.
\end{sourceguide}

\section{Why doubled angles need a separate theory}

For principal angles $\theta_j\in[0,\pi/2]$, the doubled sine and tangent are
\[
  \sin(2\theta_j)=2\sin\theta_j\cos\theta_j,
\]
\[
  \tan(2\theta_j)=\frac{2\tan\theta_j}{1-\tan^2\theta_j}.
\]
The sine is finite everywhere but is not monotone after $\pi/4$. The tangent has a pole at $\pi/4$ and changes sign across it. Therefore:
\begin{itemize}
\item an operator $\sin(2\Theta)$ theorem can be stated without an acute branch;
\item a scalar theorem involving only the maximal angle may discard information;
\item a $\tan(2\Theta)$ theorem must prove or assume that the angle spectrum stays on a chosen side of $\pi/4$.
\end{itemize}

\section{Reflection reduction}

Let $P$ project onto the unperturbed target subspace $U$, and let $Q$ project onto a comparison subspace $V$. Define the reflection
\[
  J_Q:=2Q-I.
\]
On a principal two-plane, $J_QJ_P$ is rotation by $2\theta$. Thus the single-angle geometry between $U$ and $J_QU$ is the doubled geometry between $U$ and $V$.

Let $A=A^*$ and assume $U$ reduces $A$. If $V$ reduces $B=B^*$, then $J_QBJ_Q=B$. Hence
\[
  A-J_QAJ_Q
  =(A-B)-J_Q(A-B)J_Q.
\]
Every UI norm satisfies
\[
  N(A-J_QAJ_Q)\le2N(A-B).
\]
Applying a $\sin\Theta$ theorem to $A$ and $J_QAJ_Q$ yields the doubled-sine perturbation theorem.

This is the clean parent proof. A direct expansion in principal coordinates should be reserved for the geometric identification of the cross block with $\sin(2\Theta)$.

\section{Classical ordered finite-gap $\sin 2\Theta$}

Davis--Kahan impose a geometric separation stronger than arbitrary disjoint spectra: the convex hull of one component does not meet the other, or equivalently a relevant ordered interval/exterior configuration exists. Then the Sylvester constant is one. The reflection proof gives
\[
  \delta N(\sin(2\Theta(U,V)))
  \le2N(B-A)
\]
for arbitrary unitarily invariant norms.

This theorem is sharp in its classical configuration. It should remain distinct from the generic separated-spectrum theorem.

\section{Seelmann's generic $\sin 2\Theta$ theorem}

Let
\[
  \spec(A)=\sigma\cup\Sigma,
  \qquad
  d=\dist(\sigma,\Sigma)>0,
\]
with no convex-hull separation. Let $V$ be a bounded self-adjoint perturbation and let $Q$ project onto any reducing subspace for $A+V$. Seelmann proves
\[
  \|\sin(2\Theta(E_A(\sigma),Q))\|
  \le \frac{\pi}{2}\frac{2\|V\|}{d}
  =\pi\frac{\|V\|}{d}.
\]
The projection $Q$ need not initially be a spectral projection for $A+V$.

The proof replaces the ordered constant-one cross-block estimate by the generic $\pi/2$ spectral-subspace estimate. In symmetrically normed ideals the same architecture gives the corresponding ideal-gauge statement.

\subsection{Operator theorem versus maximal-angle theorem}

Seelmann emphasizes
\[
  \sin(2\|\Theta\|)\le\|\sin(2\Theta)\|.
\]
Equality holds when the angle spectrum is contained in $[0,\pi/4]$, but not in general. A maximal-angle scalar theorem cannot serve as the parent for arbitrary UI norms or even for the strongest operator-norm statement.

This validates the repository's decision to state \code{sinTwoTheta_generalSeparation} at the operator level.

\section{Residual $\sin 2\Theta$: the signature audit}

Suppose $X:F\to E$ is isometric and $M=M^*$. A naive desired theorem is
\[
  \delta N(\sin(2\Theta(U,\range X)))
  \le2N(AX-XM).
\]
This is not valid from an arbitrary coordinate operator $M$ alone. The reflection argument needs a comparison operator that is invariant under the reflection across $\range X$. The canonical choice is the block-diagonal operator formed from the Ritz compression and its complementary compression, or equivalently a Galerkin condition ensuring the residual is off-diagonal:
\[
  X^*(AX-XM)=0.
\]

Thus a correct parent residual theorem should use one of the following equivalent packages:
\begin{enumerate}
\item $M=X^*AX$ and the exact Ritz residual;
\item an explicit Galerkin hypothesis $X^*R=0$;
\item a reflected comparison operator $B$ reducing $\range X$ with $B-A$ controlled by the residual;
\item a stated block defect larger than the raw residual.
\end{enumerate}

The formalization must audit \code{sinTwoTheta_residual_le} before filling its proof. A proof that silently rewrites arbitrary $M$ as the Ritz matrix would establish a different theorem.

\section{The off-diagonal $\tan 2\Theta$ mechanism}

The strongest doubled-tangent results arise when the perturbation is off-diagonal relative to the unperturbed spectral split. In block form,
\[
  A=\begin{bmatrix}A_+&0\\0&A_-\end{bmatrix},
  \qquad
  V=\begin{bmatrix}0&B\\B^*&0\end{bmatrix}.
\]
The diagonal blocks do not move directly; the rotation is driven entirely by $B$. This structure gives both a stable spectral split and the acute branch.

Grubi\v{s}i\'c--Kostrykin--Makarov--Veseli\'c treat a more general indefinite quadratic-form setting. Under a relative off-diagonal bound $v$, their main estimate is
\[
  \|P-Q\|
  \le
  \sin\left(\frac12\arctan v\right)
  <\frac{\sqrt2}{2},
\]
where $P$ and $Q$ are positive spectral projections before and after perturbation. Equivalently, for the maximal angle $\theta$,
\[
  \tan(2\theta)\le v.
\]
The strict bound $\|P-Q\|<\sqrt2/2$ is exactly the branch conclusion $\theta<\pi/4$.

In bounded finite dimensions the form machinery specializes to an operator block estimate. The reusable content is:
\begin{enumerate}
\item off-diagonal structure;
\item a sign or reflection operator $J$;
\item positivity of the shifted indefinite form;
\item sectorial control of the polar unitary;
\item conversion from the polar-unitary angle to projection distance.
\end{enumerate}

\section{$\tan 2\Theta$ as a proof-carrying operator}

On the branch where $\cos(2\Theta)$ is invertible and positive, define
\[
  \tan(2\Theta)
  :=\sin(2\Theta)\cos(2\Theta)^{-1}.
\]
The definition should carry a proof that
\[
  \spec(\Theta)\subset[0,\pi/4).
\]
A totalized function assigning an arbitrary value at the pole is inappropriate for the analytic parent theorem because it can make a false statement appear total.

The branch proof may come from:
\begin{itemize}
\item off-diagonal perturbation structure;
\item a half-gap smallness theorem;
\item continuation from $t=0$ together with a spectral exclusion around $\pi/4$;
\item a direct projection bound $\|P-Q\|<\sqrt2/2$.
\end{itemize}

\section{Canonical spectral selection by continuation}

Let
\[
  A_t=A+tH,
  \qquad t\in[0,1].
\]
Choose a contour $\Gamma$ or a fixed interval/exterior decomposition separating the target cluster from its complement along the entire path. Define
\[
  P_t=\frac{1}{2\pi i}\int_\Gamma(z-A_t)^{-1}\,dz.
\]
In finite dimension, resolvent continuity implies norm continuity of $P_t$. Rank is integer-valued and locally constant, hence constant. If $\|P_t-P_s\|<1$ for nearby parameters, their ranges are graph-related and lie in the same projection component.

The endpoint $P_1$ is thereby the unique spectral branch continued from $P_0$. This is the geometric content needed for Davis--Kahan Theorem 8.1 and for honest canonical $\tan 2\Theta$ statements.

A theorem that defines the perturbed subspace only at $t=1$ using a numerical Borel set may choose the wrong component if eigenvalue neighborhoods overlap or if both reducing branches satisfy a coarse separation property.

\section{Repulsion and the acute branch}

Theorem 8.1 contains more than projection continuity. It proves that the selected eigenvalues are repelled from the gap boundary. In principal coordinates, the cosine block compares compressions of $A$ and $B$. Positivity and min--max arguments show that eigenvalues on the selected branch cannot cross inward through the gap.

Theorem 8.2 supplies a quantitative criterion, often a half-gap condition, ensuring the selected projection remains within norm one of the original projection. This gives a graph representation. A stronger bound below $\sqrt2/2$ selects the $\pi/4$ branch needed by doubled tangent.

The formal dependency should be:
\[
  \text{spectral inclusion/repulsion}
  \Longrightarrow
  \text{projection component}
  \Longrightarrow
  \text{acute or quarter-angle branch}
  \Longrightarrow
  \tan(2\Theta)\text{ theorem}.
\]
Do not reverse this by assuming an unnamed ``canonical'' acute subspace in the tangent theorem.

\section{Maximal finite theorem matrix}

\begin{longtable}{@{}p{0.20\textwidth}p{0.31\textwidth}p{0.40\textwidth}@{}}
\toprule Configuration & Correct parent result & Constant and branch \\
\midrule\endfirsthead
\toprule Configuration & Correct parent result & Constant and branch \\
\midrule\endhead
Ordered/convex-hull separation & Every-UI $\sin 2\Theta$ perturbation theorem & $2/\delta$; no acute branch needed. \\
Arbitrary separated spectra & Operator/ideal $\sin 2\Theta$ via generic Sylvester & $\pi/\delta$; operator theorem stronger than scalar maximal-angle form. \\
Ritz residual & Every-UI residual $\sin 2\Theta$ & Requires Galerkin/off-diagonal residual structure. \\
Off-diagonal perturbation & $\tan 2\Theta$ and projection-distance bound & Sharp half-arctangent form; branch $\Theta<\pi/4$ is derived. \\
General perturbation, canonical cluster & Branch-selected $\tan 2\Theta$ & Needs continuation and repulsion; not valid for arbitrary reducing subspace. \\
\bottomrule
\end{longtable}

\section{Lean-sized proof seams}

\begin{enumerate}
\item Centralize reflection and doubled-angle identities.
\item Prove the UI reflection-defect bound.
\item Audit the residual theorem and add the minimal Galerkin hypothesis.
\item Prove the generic $\pi$ operator-valued theorem from the sharp generic Sylvester result.
\item Define a proof-carrying quarter-angle branch object.
\item Prove continuous spectral projections along $A_t$.
\item Prove projection-component and range-equivalence lemmas from norm $<1$.
\item Prove spectral repulsion and endpoint identification.
\item Derive the canonical finite $\tan 2\Theta$ theorem.
\item Add planar sharpness models on both sides of the branch threshold.
\end{enumerate}

\section{Repository map}

\begin{formalizationmap}
Reflection defect & Reduction of doubled sine to a single-angle theorem & \code{SinTwoTheta.lean} has a green perturbation route; residual parent still needs audit. \\
Generic separation & $\pi$ constant and symmetrically normed ideals & Experimental \code{DoubleAngle.lean}; public finite theorem should follow once generic Sylvester is unconditional. \\
Quarter-angle branch & Proof-carrying $\tan(2\Theta)$ & \code{TanTwoTheta.lean}; several branch and selection results remain sorry. \\
Continuation & Norm-continuous Riesz/spectral projection path & Experimental \code{Continuation.lean}; primary missing topological root. \\
Repulsion & Theorem 8.1 compression and eigenvalue movement & Signatures present in \code{TanTwoTheta.lean}; proofs incomplete. \\
Half-gap criterion & Theorem 8.2 acute selection & \code{Generalized.lean} and \code{TanTwoTheta.lean}; should be downstream of continuation. \\
\end{formalizationmap}

\section*{Bibliography}
\begin{thebibliography}{9}
\bibitem{DK1970} C. Davis and W. M. Kahan, \emph{The Rotation of Eigenvectors by a Perturbation. III}, SIAM J. Numer. Anal. 7 (1970), 1--46.
\bibitem{Seelmann2013} A. Seelmann, \emph{Notes on the $\sin 2\Theta$ Theorem}, arXiv:1310.2036.
\bibitem{GKMV2010} L. Grubi\v{s}i\'c, V. Kostrykin, K. A. Makarov, and K. Veseli\'c, \emph{The $\tan 2\Theta$ Theorem for Indefinite Quadratic Forms}, arXiv:1006.3190.
\bibitem{AM2013} S. Albeverio and A. K. Motovilov, generic $\sin 2\theta$ estimates for spectral subspaces, Complex Analysis and Operator Theory 7 (2013), 1389--1416.
\end{thebibliography}

\end{document}
